\chapter{Cơ sở lý thuyết}

\section{Mô hình ngôn ngữ lớn}

\subsection{Kiến trúc Transformer}

Sự ra đời của kiến trúc Transformer đã đặt nền móng cho sự bùng nổ của trí tuệ nhân tạo và lĩnh vực xử lý ngôn ngữ tự nhiên (NLP) hiện đại. Mô hình này là một bước ngoặc lớn trong việc xử lý dữ liệu dạng chuỗi. Khác với Transformer, trong quá khứ, các mô hình xử lý ngôn ngữ chủ yếu dựa trên kỹ thuật N-grams kết hợp với các mô hình xác suất toán học, vốn thường có khả năng dự đoán bị giới hạn nghiêm trọng do chỉ dựa vào bối cảnh dữ liệu hạn hẹp của từ vựng. Chính vì thế, các mô hình học sâu (Deep Learning) đã được phát triển để giải quyết vấn đề này. Có thể kể đến các mô hình đại diện nổi bật như Mạng nơ-ron hồi quy (RNN) và Bộ nhớ ngắn hạn dài (LSTM), cho phép xử lý chuỗi văn bản một cách tuần tự, nhờ đó giải quyết vấn đề ngữ cảnh. Tuy nhiên, các mô hình có kiến trúc này đối mặt hai vấn đề kỹ thuật nghiêm trọng. Thứ nhất, chúng không cho phép tính toán dữ liệu một cách song song, làm tăng đáng kể thời gian huấn luyện và không tận dụng được sức mạnh phần cứng. Thứ hai, các mô hình này gặp vấn đề vanishing gradient khi phải xử lý các chuỗi văn bản quá dài.

Để giải quyết bài toán trên, Transformer đã loại bỏ hoàn toàn cơ chế tính toán tuần tự và thay thế bằng cơ chế tự chú ý (Self-Attention). Cơ chế này tính toán và đánh giá mức độ tương quan giữa các token trong văn bản đầu vào cùng một lúc thông qua công thức toán học sau:

\[
\text{Attention}(Q, K, V) = \text{softmax}\left(\frac{QK^T}{\sqrt{d_k}}\right)V
\]

Trong phương trình trên, các ma trận $Q$ (Query), $K$ (Key) và $V$ (Value) là các vector tính toàn từ văn bản đầu vào. Tham số $d_k$ đại diện cho số chiều của vector $K$. Ngoài ra, để bổ sung thông tin về vị trí của mỗi token trong các vector, Transformer tích hợp thêm cơ chế Mã hóa vị trí (Positional Encoding), nó cung cấp cho mỗi token một giá trị vị trí riêng biệt.

Tận dụng nền tảng mô hình mạnh mẽ này của Transformer, các mô hình Ngôn ngữ Lớn ngày này đang phát triển không ngừng. Các hệ thống Mô hình Ngôn ngữ Lớn có thể được chia thành ba nhánh kiến trúc chính. Nhóm đầu tiên là các mô hình chỉ có Bộ mã hóa (Encoder-only) như họ BERT, các mô hình này có điểm mạnh là hiểu ngữ cảnh do cơ chế self-attention hai chiều nhưng lại tỏ ra kém hiệu quả trong việc sinh văn bản. Nhóm thứ hai có kiến trúc đầy đủ gồm 1 bộ Mã hóa và một bộ Giải mã (Encoder-Decoder) như T5 hay BART, các mô hình này phù hợp cho các tác vụ thiên về chuyển đổi chuỗi dữ liệu sang một chuỗi dữ liệu khác như các bài toán dịch máy. Nhóm cuối cùng là nhóm các mô hình chỉ có bộ Giải mã (Decoder-only), đây chính là nhóm kiến trúc đang thống trị và được nghiên cứu nhiều nhất hiện nay. Các đại diện có thể kể đến của dòng mô hình này như GPT, Gemini, Qwen, LLaMA. Đặc điểm nổi bật của kiến trúc này là chúng được tối ưu cho tác vụ dự đoán token tiếp theo theo phương pháp tự hồi quy (autoregressive generation). Nhờ vào lượng dữ liệu khổng lồ và sự gia tăng số lượng tham số lên đến quy mô hàn trăm chục, trăm tỷ, đã tạo ra các LLMs có khả năng suy luận logic phức tạp, hiểu sâu ngữ cảnh chưa từng có.


\subsection{Các mô hình ngôn ngữ lớn sử dụng trong nghiên cứu}

Trong bối cảnh hệ sinh thái AI đang phân hóa và phát triển với tốc độ chóng mặt,nghiên cứu này sẽ tập trung đánh giá trên hai nhóm mô hình ngôn ngữ lớn: nhóm mã nguồn mở (Open-source LLMs) và nhóm thương mại mã nguồn đóng (Proprietary/Closed-source LLMs).

\textbf{Nhóm mô hình mã nguồn mở:}

Sự phát triển nhanh chóng của các mô hình mã nguồn mở đang dần thu hẹp khoảng cách giữa các tập đoàn công nghệ lớn và các cộng đồng học thuật, công ty công nghệ nhỏ đến trung bình. Trong nghiên cứu này, ba mô hình đại diện tiêu biểu sau được lựa chọn để đánh giá:

\begin{itemize}
    \item \textbf{Gemma 3 12B:} Mô hình này được Google phát hành chính thức vào tháng 3 năm 2025, thuộc họ kiến trúc đa phương thức với tổng cộng 12.2 tỷ tham số. Mục tiêu của Gemma 3 12B là có thể chạy trực tiếp trên các thiết bị máy tính cá nhân hoặc hạ tầng máy chủ tại chỗ. Ưu điểm nổi bật của mô hình này là kích thước cửa sổ ngữ cảnh rất lớn lên tới 128,000 token đối với văn bản. Nhờ đặc điểm này mà Gemma 3 có thể xử lý được các các tài liệu, giáo trình phức tạp mà không bị cắt xén hay mất mát thông tin quan trọng. 
    \item \textbf{Qwen 3 14B:} Là một trong các mô hình thuộc dòng Qwen nổi tiếng của Alibaba Cloud, được ra mắt vào tháng 4 năm 2025 với 14.8 tỷ tham số.  Qwen 3 14B đã thể hiện được sự vượt trội trong các bài kiểm tra về khoa học tự nhiên, toán học và lập trình. Mặc dù Qwen 3 có cửa sổ ngữ cảnh khá hạn chế, bị giới hạn ở mức 33,000 đến 41,000 token (hẹp hơn đáng kể so với Gemma 3), nhưng bù lại, nó cung cấp khả năng gọi hàm (function calling) cực tốt và xử lý định dạng JSON rất chính xác và ổn định.
    \item \textbf{Phi 4:} Được Microsoft phát triển và công bố vào tháng 12 năm 2024, Phi 4 được huấn luyện trên lượng dữ liệu cực kỳ chất lượng thay vì sử dụng dữ liệu hỗn tạp từ internet. Với kiến trúc 14 tỷ tham số với cửa sổ ngữ cảnh 16,000 token, mô hình được huấn luyện trên lượng dữ liệu khổng lồ 9.8 nghìn tỷ token chất lượng cao từ các sách học thuật chuyên sâu và các bộ dữ liệu hỏi đáp QA đã qua tinh lọc. Chính vì triết lý đặc biệt này, Phi 4 đã đạt hiệu suất đáng kinh ngạc trong các bài đánh giá STEM, thậm chí vượt qua cả hiệu suất của các mô hình khổng lồ như GPT-4 trong một số lĩnh vực hỏi đáp chuyên sâu.
\end{itemize}

\textbf{Nhóm mô hình thương mại:}

Để có thể đánh giá hiệu suất của các mô hình ngôn ngữ mã nguồn mở trên trong bài toán sinh câu hỏi trắc nghiệm, nghiên cứu đã sử dụng các mô hình thương mại sau để thiết lập một bộ đánh giá khách quan và đáng tin cậy:

\begin{itemize}
    \item \textbf{Gemini 3.0 Flash:} Thuộc thế hệ mô hình đa phương thức được Google phát hành vào tháng 12 năm 2025, Gemini 3.0 Flash được thiết kế với mục tiêu phản hồi tức thì. Không những thế, với cửa sổ ngữ cảnh lên đến 1 triệu token, mô hình cung cấp một khả năng có thể xử lý lượng dữ liệu đầu vào khổng lồ mà không gặp hiện tượng suy giảm trí nhớ giữa chừng. Mặc dù tỷ lệ ảo giác của mô hình có thể cao, sự bù đắp đến từ độ chính xác 88\% và khả năng tự phân tích lỗi là một đánh đổi chấp nhận được.
    \item \textbf{GPT 5 mini:} Được phát hành vào tháng 8 năm 2025 bởi OpenAI thuộc dòng mô hình cân bằng giữa tốc độ trả lời và hiệu năng. GPT 5 mini có cửa sổ ngữ cảnh lên đến 400,000 token và hỗ trợ đầu vào đa phương thức. Với mức giá hợp lý cùng khả năng xuất dữ liệu có cấu trúc JSON chính xác và ổn định, GPT 5 mini là một mô hình đáng tin cậy để sử dụng cho bài toán tạo câu hỏi trắc nghiệm.
    \item \textbf{Grok 4.1 fast:} Tương tự GPT 5 mini, mô hình này thuộc dòng mô hình tốc độ cao do xAI phát hành vào tháng 11 năm 2025. Mô hình được đánh giá cực kỳ cao nhờ tỷ lệ ảo giác cực thấp (chỉ 4\% dựa trên bài kiểm tra FActScore). Nhờ khả năng phân tích đa phương thức hiệu quả, tính năng sử dụng các công cụ bên ngoài và sự tích hợp sâu vào nền tảng X, Grok 4,1 fast có thể cung cấp những câu trả lời có độ chính xác cao.
\end{itemize}

\begin{table}[htbp]
    \centering
    \caption{Thông số kỹ thuật các mô hình ngôn ngữ lớn}
    \label{tab:llm_comparison_fit}
    \renewcommand{\arraystretch}{1.5}
    \resizebox{\textwidth}{!}{%
        \begin{tabular}{|l|l|c|c|c|l|}
            \hline
            \textbf{Mô hình} & \textbf{Nhà phát triển} & \textbf{Số tham số} & \textbf{Cửa sổ ngữ cảnh} & \textbf{Khả năng thị giác} & \textbf{Ngày phát hành} \\
            \hline
            Gemma-3-12b-it & Google & 12.2B & 128K & Có & Tháng 3, 2025 \\
            \hline
            Qwen3-14b & Alibaba Cloud & 14.8B & 33K - 41K & Không & Tháng 4, 2025 \\
            \hline
            Phi-4 & Microsoft & 14B & 16K & Không & Tháng 12, 2024 \\
            \hline
            Gemini-3.0-flash & Google & - & 1,000K & Có & Tháng 12, 2025 \\
            \hline
            GPT-5-mini & OpenAI & - & 400K & Có & Tháng 8, 2025 \\
            \hline
            Grok-4.1-fast & xAI & - & 2,000K & Có & Tháng 11, 2025 \\
            \hline
        \end{tabular}%
    }
\end{table}

\section{Sinh câu hỏi trắc nghiệm tự động}

Bài toán Tạo câu hỏi trắc nghiệm tự động (Automated Multiple-Choice Question Generation - AMCQG) đã chuyển mình từ các phương pháp sử dụng các quy tắc ngôn ngữ sơ khai, quy luật sang việc ứng dụng các mô hình AI tạo sinh hiện đại, nổi bật là các LLMs. Đây là cột mốc quan trọng đánh dấu sự chuyển dịch tất yếu từ việc kiểm tra khả năng ghi nhớ sang việc đánh giá năng lực phân tích và giải quyết vấn đề của người học.

\subsection{Các phương pháp truyền thống}

Trước thời điểm bùng nổ của Mô hình các ngôn ngữ lớn, bài toán Sinh câu hỏi tự động thường được giải quyết bằng các phương pháp Xử lý Ngôn ngữ Tự nhiên truyền thống. Đây là các phương pháp tiếp cận mang thiên hướng thống kê và ngôn ngữ học. Quy trình AMCQG thường được chia thành ba bước xử lý tuần tự: trích xuất khái niệm, luận điểm cốt lõi, sinh câu hỏi và xây dựng các đáp án gây nhiễu.

Lý thuyết cốt lõi đằng sau các hệ thống AMCQG truyền thống là dựa vào các quy luật ngôn ngữ (Rule-based) và sử dụng các mẫu có sẵn (Template-based). Các hệ thống này sẽ phân tích cú pháp của đoạn văn bản ban đầu thông qua các kỹ thuật NLP như Nhận dạng thực thể có tên (Named Entity Recognition) và Gán nhãn từ loại (art-of-Speech Tagging). Ví dụ, khi văn bản có chứa một chủ ngữ là “Nikola Tesla” và một cụm từ chỉ địa điểm như “thành phố New York”, hệ thống sẽ tự động nhận diện và lựa chọn một mẫu câu hỏi có sẵn như (Wh) + Động từ + Chủ thể để tạo ra một câu hỏi như "Nikola Tesla hiện đang làm việc ở thành phố nào?”. Tiếp theo, hệ thống sẽ thực hiện sinh ra các đáp án nhiễu. Mô hình có thể sử dụng các kỹ thuật khai thác kho từ vựng phân cấp (Ontology), hoặc bộ dữ liệu từ vựng ngữ nghĩa (WordNet, Word2Vec) để tìm kiếm các từ đồng nghĩa hoặc các thực thể có chung đặc tính, tính chất (ví dụ như tìm cá thành phố khác ngoài New York để làm đáp án sai).

Mặc dù các phương pháp tạo câu hỏi trắc nghiệm truyền thống sử dụng kỹ thuật NLP này đem lại khả năng kiểm soát dữ liệu chặt chẽ kèm theo tính minh bạch cao và tính đúng đắn tuyệt đối về mặt ngữ pháp, chúng lại tồn tại những điểm yếu chí mạng khiến chúng chưa thực sự đem lại hiệu quả cao cũng như giải quyết được các nhu cầu sư phạm thực tế. Thứ nhất, hệ thống này yêu cầu một lượng tài nguyên về con người vô cùng lớn; chúng yêu cầu sự tinh chỉnh chính xác từ các chuyên gia ngôn ngữ học để có thể xây dựng và bảo trì hàng nghìn quy tắc khác nhau cho từng ngôn ngữ và chủ đề. Thứ hai, phương pháp này chỉ sinh ra các câu hỏi đơn giản dạng truy xuất dữ liệu bề mặt hoặc điền vào chỗ trống, các câu hỏi này không yêu cầu người đọc phải liên kết thông tin giữa nhiều đoạn văn của bài đọc dẫn đến các câu hỏi đầu ra thiếu hoàn toàn tính chất tư duy và suy luận sâu. Thứ ba, chất lượng của các đáp án sai sinh ra rất kém và thiếu tính gẫy nhiễu; chúng dễ dàng bị người học loại trừ bằng các mẹo cơ bản vì thiếu tính liên kết logic, đồng thời không phản ánh được những lỗi sai phổ biến mà học sinh thường mắc phải trong quá trình giải bài. Những hạn chế kể trên là nguyên nhân chính khiến cho các hệ thống NLP truyền thống không thể đáp ứng các tiêu chuẩn giáo dục hiện đại.


\subsection{Sinh câu hỏi dựa trên mô hình ngôn ngữ lớn}

Kể từ giai đoạn năm 2022 trở đi, ngành Trí tuệ nhân tạo đã chứng kiến một sự chuyển dịch triệt để trong công nghệ giáo dục nhờ sự phát triển đột phá của các Mô hình Ngôn ngữ Lớn. Khác với cách tiếp cận thống kê và quy tắc ngôn ngữ như các mô hình truyền thống. LLMs cho phép giải quyết các bài toán AMCQG một cách dễ dàng, trong đó các mô hình AI không chỉ thao tác cơ học trên các từ ngữ mà chúng thực sự hiểu văn bản về mặt ngữ nghĩa.

Các nhà giáo dụng thường sử dụng kỹ thuật Prompt để có thể yêu cầ LLM phân tích văn bản đầu vào và sinh ra các câu hỏi theo yêu cầu cùng với các đáp án nhiễu. Đối với các hệ thống tự động, LLM thường được yêu cầu để sinh ra dữ liệu có định dạng JSON hoặc XML. Những nghiên cứu gần đây còn tích hợp các kỹ thuật cao cấp như Sinh tăng cường truy xuất (Retrieval-Augmented Generation - RAG) và Mạng sinh đối nghịch (Generative Adversarial Networks - GANs). Đối với mô hình RAG, một cơ sở dữ liệu tri thức khổng lồ sẽ được sử dụng; trước khi sinh câu hỏi, các LLM sẽ truy vấn các tài liệu liên quan từ cơ sở dữ liệu này một cách chính xác nhất thông qua chỉ mục, điều này giảm thiểu đáng kể hiện tượng ảo giác và đảm bảo các câu hỏi sinh ra luôn được gắn liền với một sự thật có thể kiểm chứng.

Ngoài ra, sự vượt trội của LLMs so với các kỹ thuật truyền thống thể hiện ở khả năng thiết kế ra các phương án nhiễu mang tính chất chẩn đoán cáo. Thay vì chỉ đơn giản lựa chọn các từ đồng nghĩa, một LLM khi được tinh chỉnh tốt có thể chủ động mô phỏng một cách chính xác các chuỗi suy luận, tư duy bị lỗi của người học và tạo ra các đáp án sai nhưng cực kỳ hợp lý. Chính năng lực này đã đưa các bài toán AMCQG vượt ra khỏi giới hạn của các bài toán kiểm tra trí nhớ cơ học. Nói cách khác, hệ thống khi tận dụng LLMs có thể tạo ra các hỏi ở các cấp độ nhận thức cao hơn.


\section{Phân loại Bloom}

Để các hệ thống sinh câu hỏi trắc nghiệm tự động có thể vừa đảm bảo các mục tiêu sư phạm và vừa không biến chúng thành môt  cỗ máy tạo câu hỏi ngẫu nhiên, việc thiết kế độ khó cần dựa trên một thang đo đã được kiểm chứng bởi khoa học. Nghiên cứu này sẽ tập trung vào khung nhận thức khoa học được công nhận rộng rãi nhất trong lĩnh vực giáo dục: Phân loại Bloom (Bloom’s Taxonomy). Thang đo này được nhà tâm lý học và giáo dục Benjamin Bloom đề xuất lần đầu vào năm 1956 và được tinh chỉnh bởi Anderson và Krathwohl vào năm 2001. Trong nghiên cứu này, nhằm mục đích tối ưu hóa và phù hợp với hình thức thi trắc nghiệm khách quan đa lựa chọn (MCQ), sáu cấp độ ban đầu của thang đo Bloom được phân loại lại thành ba cấp độ cốt lõi: Truy xuất, Suy luận và tổng hợp và Lập luận phản biện. 

\subsection{Cấp độ 1 - Truy xuất}

Cấp độ Truy xuất tương ứng với bậc “Nhớ” (Remember) trong thang đo Bloom sửa đổi. Cấp độ này là nền tảng của sự nhận thức, tập trung chủ yếu vào khả năng ghi nhớ cơ học và truy xuất thông tin dài hạn, hoặc nhận dạng một đối tượng, sự kiện được trình bày rõ ràng trong một đoạn văn bản thuộc bài đọc. Quá trình này thường được gọi là quá trình nhận thức nông.

Trong quá trình tinh chỉnh các chỉ dẫn cho LLMs sinh câu hỏi ở cấp độ Truy xuất, hệ thống được cài đặt để tập trung vào các chi tiết cụ thể, định nghĩa thuật ngữ và nhận dạng các sự kiện. Đặc thù của của MCQ ở dạng này thường xoay quanh các câu hỏi “Ai?”, “Cái gì?”, “Ở đâu?” và “Khi nào?”. Đối với bài toán  này, LLMs thể hiện vô cùng xuất sắc trong việc tự động tạo ra các câu hỏi, bởi vì quy trình này cơ bản chỉ đòi hỏi khả năng đối sánh chuỗi và trích xuất các thực thể từ văn bản mà không cần hiểu sâu hay xây dựng bất kỳ chuỗi suy luận logic nào. Không chỉ vậy, các phương án nhiễu ở mức độ này cũng tương đối dễ xây dựng; chúng thường là các thực thể hoặc dữ liệu xuất hiện trong văn bản gốc để tạo ra sự khó khăn cho người đọc.

\subsection{Cấp độ 2 - Suy luận và tổng hợp}

Cấp bậc thứ hai bao gồm hai kỹ năng Suy luận và Tổng hợp, đại diện cho các bậc “Hiểu” (Understand), “Áp dụng” (Apply) và “Sáng tạo” (Create) trong phân loại Bloom. Ở cấp độ này, thay vì được hỏi các câu hỏi có thể truy xuất trực tiếp từ bề mặt của văn bản, người học sẽ buộc phải xử lý thông tin đã học và ứng dụng lý thuyết đó vào một ngữ cảnh hoàn toàn mới hoặc một tình huống giả định.

Trong bối cảnh bài toán MCQ, Suy luận yêu cầu người học phải có khả năng liên kết các mệnh đề nằm rải rác trong bài đọc để rút ra các kết quả, nhận định không được tác giả viết ra một cách trực tiếp. Đối với Tổng hợp, cấp độ này yêu cầu người học phải chắt lọc, sắp xếp lại các yếu tổ riêng lẻ thành một góc nhìn bao quát hoặc dự đoán một kết quả tương lai. Bài toán này đòi hỏi cần có các chỉ dẫn cụ thể và hợp lý cho LLMs. Ngoài ra, việc tạo ra các đáp án nhiễu ở cấp độ này cũng là một thách thức kỹ thuật đối với AI: các đáp án sai phải được thiết kế dựa vào sự suy diễn sai lầm về mặt logic hoặc dựa vào việc người đọc hiểu sai một quy luật nào đó để từ đó có thể tạo ra các lựa chọn bẫy hợp lý về mặt trực giác. 

\subsection{Cấp độ 3 - Lập luận phản biện}

Cấp độ cuối cùng của nghiên cứu này là Cấp độ Lập luận phản biện, tương ứng vào bậc “Phân tích” (Analyze) và “Đánh giá” (Evaluate). Trọng tâm của cấp độ nhận thức này là khả năng chia rỏ một lượng thông tin đồ sộ, phức tạp thành các thành phần nhỏ hơn nhằm phân tích, nhận diện các giả định ngầm, và đưa ra những nhận xét, quan điểm khách quan dựa trên sự dối chiếu với các bằng chứng thực tiễn.

Các câu hỏi trắc nghiệm ở cấp độ này thường yêu cầu người học tưởng tượng ra các tình huống kiểm tra ngụy biện logic, yêu cầu họ phân biệt rõ ràng giữa một sự thật khách quan và một quan điểm chủ quan, hoặc được yêu cầu đánh giá độ chính xác, hợp lý của một lập luận khoa học. Thách thức lớn nhất mà LLMs phải đối mặt trong cấp độ này việc thiết kế ra các câu hỏi với chất lượng cao cùng một lựa chọn đáp án đúng duy nhất. Quan trọng hơn cả, phương án nhiễu ở mức độ này không thể chỉ là những thông tin sai lệch thông thường, mà phải là các lập luận, lý lẽ nhìn bề mặt có vẻ cực kỳ thuyết phục nhưng lại chưa các lỗ hổng tinh tế về mặt logic, từ đó buộc người học phải sử dụng tư duy phản biệt cao đội mới có thể nhận diện được.

\section{Hệ thống đa tác tử}

\subsection{Khái niệm hệ thống đa tác tử}

Hệ thống đa tác tử (Multi-Agent Systems) đại diện cho một thế hệ kiến trúc Trí tuệ Nhân tạo hiện đại, cho phép giải quyết các bài toán phức tạp mà một LLM duy nhất không thể xử lý bằng cách chia nhỏ bài toàn thành các nhiệm vụ nhỏ hơn và giao cho một mạng lưới gồm nhiều tác nhân chuyên biệt, có khả năng tương tác qua lại. Một sai lầm phổ biến của nhiều hệ thống là yêu cầu một LLM giải quyết một bài toán đòi hỏi quy trình phức tạp về logic và nghiệp vụ khiến kết quả đầu ra các chất lượng kém và không thể sử dụng. Nguyên nhân chính là việc một LLM bị giao quá nhiều vai trò và nhiệm vụ như vừa phải đọc hiểu hàng nghìn từ, sinh câu hỏi, tạo ra đáp án nhiễu vào trong một lần gọi LLM duy nhất.

Mô hình Multi-Agent Systems giải quyết triệt để những vấn đề này bằng cách thiết kế một hệ thống mô phỏng cách vận hành của con người và dựa trên triết lý nền tảng của ReACT (Reasoning và Acting). Khác với việc việc trả lời một cách thụ động, mô hình ReACT buộc LLMs phải thực hiện một vòng lặp nội tại: Suy nghĩ, Hành động, Quan sát kết quả, Điều chỉnh suy nghĩ. Trong mô hình Multi-Agent Systems, mỗi tác tử được trang bị một một chỉ thị cực kỳ hẹp nhưng có tính chất chuyên môn hóa cực cao, cho phép mỗi tác tử có thể tập trung toàn bộ tài nguyên nhận thức vào một vấn đề duy nhất. Ví dụ, Tác tử Phân tích chỉ chịu trách nghiệm lọc và cấu trúc hóa văn bản thô; sau đó kết quả được chuyển cho các Tác tử Sinh câu hỏi để tạo ra bộ câu hỏi ban đầu, kết quả tiếp tục được chuyển cho cho hai Tác tử Đánh giá và Sửa lỗi để thực hiện vòng cải thiện bộ câu hỏi. Sự chuyên môn hóa này không chỉ giảm thiểu ảo giác mà còn giúp cải thiện chất lượng của nội dung sinh ra, giảm thiểu tối đa lỗi và duy trì sự ổn định cho hệ thống.

\subsection{Framework LangGraph}

Để có thể cài đặt và triển khai hệ thống đa tác tử một cách đúng đắn, nghiên cứu này sử dụng LangGraph - một thư viện mã nguồn mở chuyên dụng được xây dựng trên nền tảng của hệ sinh thái LangChain. Khác với triết lý xử lý tuần tự tuyến tính của Langchain, LangGraph được thiết kế riêng biệt cho việc tạo ra các luồng tác tử thông qua đồ thị có hướng kèm theo một đối tượng trạng thái (state) được sử dụng để lưu trữ dữ liệu khi hệ thống duyệt qua từng nút của đồ thị.

Ba thành phần chính của một hệ thống Multi-Agent Systems khi được triển khai bằng LangGraph bao gồm:

\begin{enumerate}

    \item \textbf{Trạng thái: }Là một cấu trúc dữ liệu trung tâm, nó hoạt động như một bộ nhớ chung để duy trì ngữ cảnh cho toàn bộ hệ thống. Ví dụ như các ý chính đã phân tích từ văn bản gốc, các câu hỏi đã tạo, lỗi của các câu hỏi, …
    
    \item \textbf{Nút:} Đây là các hàm Python có nhiệm vụ xử lý logic cho từng tác tử (Nút sinh câu hỏi cho cấp độ 3, Nút thực hiện đánh giá câu hỏi cấp độ 2, …). Nút này sẽ nhận được Trạng thái hiện tại, gửi lệnh đến LLM để xử lý và sau đó cập nhật Trạng thái mới với kết quả vừa tạo ra.
    
    \item \textbf{Cạnh:} Thành phần này có nhiệm vụ xác định hướng du chuyển của luồng dữ liệu giữa các Nút trong đồ thị. LangGraph cung cấp khả năng cho phép các hệ thống tự đưa ra quyết định rẽ nhánh, kích hoạt cơ chế thực thi song song (fan-out/fan-in), hoặc kết thúc luồng hoạt động dựa trên các tiêu chí xuất hiện ngay trong thời gian thực.

\end{enumerate}

Khả năng hỗ trợ cấu trúc đồ thị của LangGraph đã cho phép triển khải các hệ thống AI tự định hướng. Nếu một tác tử Đánh giá phát hiện câu hỏi sinh ra có chứa phương án nhiễu vô lý hay chứa nhiều đáp án đúng, hệ thống hoàn toàn có thể yêu cầu một Tác tử Sửa lỗi xử lý vấn đề rồi tiếp tục đưa kết quả đó cho Tác tử Sửa lỗi để đánh giá lại, chu trình này chỉ dừng khi hệ thống tạo ra các kết quả có chất lượng tốt. Ngoài ra, LangGraph còn được trang bị bộ nhớ chốt chặn giúp hệ thống có thể phục hồi trạng thái khi gặp các vấn đề như lỗi API. Đặc biệt, LangGraph tuân thủ triết lý không sử dụng prompt ẩn, mang lại cho các nhà nghiên cứu quyền kiểm soát tuyệt đối trong việc tùy chỉnh prompt và minh bạch hóa hệ thống.


\section{Nhận dạng ký tự quang học}
\subsection{OCR truyền thống}

Các mô hình OCR truyền thống bao gồm việc chuyển đổi hình ảnh chứa văn bản thành định dạng kỹ thuật số có thể thao tác được thông qua một quy trình xử lý phức tạp và độc lập. Cấu trúc của các hệ thống OCR truyền thống thường bao gồm các bước sau:

\begin{enumerate}
    \item \textbf{Tiền xử lý hình ảnh:} Hình ảnh đầu vào được cắt xén, loại bỏ nhiễu và làm biến dạng để đảm bảo chất lượng trong các giai đoạn xử lý tiếp theo.

    \item \textbf{Phân tích bố cục:} Hệ thống phân đoạn hình ảnh đã xử lý thành các thành phần của tài liệu, chẳng hạn như khối văn bản, bảng, hình ảnh và tiêu đề, bằng cách sử dụng các thuật toán dựa trên quy tắc hoặc các mô hình CNN học sâu.

    \item \textbf{Phát hiện văn bản:} Mô hình xác định vị trí các khung bao quanh chứa văn bản trong hình ảnh.

    \item \textbf{Nhận dạng văn bản:} Các vùng văn bản được phát hiện trong bước trước được chuyển đến một mô hình học sâu, chẳng hạn như Mạng thần kinh hồi quy tích chập (CNN), để trích xuất các ký tự từ các vùng này.

    \item \textbf{Xử lý hậu kỳ:} Hệ thống kiểm tra chính tả của văn bản được trích xuất và lắp ráp lại dựa trên từ điển hoặc các mô hình thống kê của ngôn ngữ.

\end{enumerate}

Mặc dù các mô hình này đạt độ chính xác cao đối với văn bản in có cấu trúc rõ ràng, OCR truyền thống có nhiều hạn chế khi xử lý các tài liệu phức tạp (như báo cáo tài chính hoặc tài liệu toán học) hoặc hình ảnh văn bản bị nhiễu và mờ. Điều này dẫn đến kết quả không chính xác và không đảm bảo thứ tự đọc chính xác.


\subsection{OCR dựa trên mô hình thị giác--ngôn ngữ}

Không giống như kiến trúc OCR đa mô-đun truyền thống, OCR dựa trên Mô hình Ngôn ngữ Thị giác (VLM) giải quyết bài toán trong một lần gọi duy nhất. VLM mã hóa trực tiếp hình ảnh đầu vào và kết hợp nó với các chỉ dẫn để tạo ra văn bản định dạng chuẩn (như Markdown hoặc JSON). Phương pháp này hoàn toàn giải quyết vấn đề thứ tự đọc và bảo toàn cấu trúc logic của tài liệu gốc.

Gần đây, sự xuất hiện của các mô hình ngôn ngữ thị giác quy mô nhỏ (thường có ít hơn 2 tỷ tham số) đã tối ưu hóa chi phí tính toán trong khi đạt được hiệu suất tương đương với các mô hình lớn hơn. Sau đây là phân tích kỹ thuật của bốn kiến trúc OCR VLM nhỏ gọn tiên tiến:

% -

\begin{itemize}
    \item \textbf{DeepSeek OCR-2:}
    \begin{itemize}
        \item \textbf{Kiến trúc:} DeepSeek OCR-2 duy trì cấu trúc Encoder-Decoder nhưng nâng cấp bộ mã hóa hình ảnh lên DeepEncoder V2. Mô hình áp dụng cơ chế Visual Causal Flow, cho phép học tuyến tính và ngữ nghĩa về thứ tự đọc của con người thay vì dựa vào các sắp xếp tọa độ hình học cố định.
        \item \textbf{Đặc điểm:} Được tối ưu hóa cao cho việc trích xuất tài liệu theo thời gian thực và xây dựng các luồng RAG. Phiên bản DeepSeek-VL2-Tiny (1.0B) đạt hiệu suất xuất sắc trên các tập dữ liệu đa phương thức với yêu cầu tài nguyên phần cứng cực thấp.
    \end{itemize}

    \item \textbf{dots.ocr:}
    \begin{itemize}
        \item \textbf{Kiến trúc:} Dựa trên một LLM nhỏ gọn với tổng cộng 1,7 tỷ tham số, dots.ocr tích hợp các tác vụ phân tích thiết kế và nhận dạng nội dung vào một VLM duy nhất mà không cần quy trình nhiều giai đoạn.
        \item \textbf{Đặc điểm:} Các tác vụ được điều khiển thông qua các lệnh. Mô hình thể hiện hiệu suất vượt trội trong việc phân tích tài liệu đa ngôn ngữ và xử lý hiệu quả các ngôn ngữ có tài nguyên thấp. Hiệu suất của dots.ocr trên tập dữ liệu OmniDocBench đạt mức tiên tiến nhất (SOTA), tương đương với các mô hình có số lượng tham số cao.
    \end{itemize}

    \item \textbf{MinerU 2.5:}
    \begin{itemize}
        \item \textbf{Kiến trúc:} MinerU 2.5, một mô hình VLM chuyên dụng với 1,2 tỷ tham số, sử dụng mô hình ngôn ngữ-hình ảnh tách rời. Mô hình sử dụng mã hóa độ phân giải gốc được phát triển từ Qwen2-VL-2B-Instruct và kết hợp với lõi ngôn ngữ từ Qwen2-Instruct-0.5B.
        \item \textbf{Tính năng:} Sau quá trình huấn luyện căn chỉnh đa phương thức chuyên sâu, MinerU 2.5 cung cấp hiệu suất trang cực nhanh (nhanh hơn tới 7 lần so với dots.ocr) và thể hiện khả năng tạo ra trực tiếp và chính xác các công thức toán học phức tạp và cấu trúc bảng HTML từ các tệp PDF.
    \end{itemize}

    \item \textbf{PaddleOCR-VL-1.5:}
    \begin{itemize}
        \item \textbf{Kiến trúc:} PaddleOCR-VL-1.5 là một kiến trúc siêu nhỏ gọn với chỉ 0,9 tỷ tham số. Mô hình tích hợp bộ mã hóa hình ảnh độ phân giải động kiểu NaViT cùng với mô hình ngôn ngữ ERNIE-4.5-0.3B.
        \item \textbf{Tính năng:} Đạt độ chính xác 94,5\% trong OmniDocBench v1.5. Tính năng mới chính của PaddleOCR-VL-1.5 là khả năng phát hiện đa giác, cho phép nó hoạt động ổn định ngay cả khi gặp phải các biến dạng vật lý trong thực tế, chẳng hạn như tài liệu bị cong vênh, bản quét mờ, phản chiếu hoặc ảnh chụp màn hình. Nó cũng hỗ trợ tự động ghép bảng và hợp nhất tiêu đề trên nhiều trang.
    \end{itemize}
\end{itemize}


\section{Phương pháp đánh giá chất lượng câu hỏi}
\subsection{Đánh giá bởi con người}
\subsection{LLM-as-a-Judge}
